\section{\label{sec:related}Related Work and Discussions}
There is a growing interest in improving the efficiency of NAS. Concurrent to our work are the promising ideas of using performance prediction~\citep{nas_performance_pred,nas_peephole}, using iterative search method for architectures of growing complexity~\citep{pnas}, and using hierarchical representation of architectures~\citep{hierarchical_nas}. Table~\ref{tab:cifar10_results} shows that \enas~is significantly more efficient than these other methods, in GPU hours.

\enas's design of sharing weights between architectures is inspired by the concept of weight inheritance in neural model evolution~\citep{nas_evolution,nas_reg_evolution}. Additionally, \enas's choice of representing computations using a DAG is inspired by the concept of stochastic computational graph~\citep{stochastic_computation}, which introduces nodes with stochastic outputs into a computational graph. \enas's utilizes such stochastic decisions in a network to make discrete architectural decisions that govern subsequent computations in the network, trains the decision maker, \ie~the controller, and finally harvests the decisions to derive architectures.

Closely related to \enas~is SMASH~\citep{smash_net}, which designs an architecture and then uses a hypernetwork~\citep{hyper_net} to generate its weight. Such usage of the hypernetwork in SMASH inherently restricts the weights of SMASH's child architectures to a low-rank space. This is because the hypernetwork generates weights for SMASH's child models via tensor products~\citep{hyper_net}, which suffer from a low-rank restriction as for arbitrary matrices $\mathbf{A}$ and $\mathbf{B}$, one always has the inequality: $\text{rank}(\mathbf{A} \cdot \mathbf{B}) \leq \min{\{\text{rank}(\mathbf{A}), \text{rank}(\mathbf{B})\}}$. Due to this limit, SMASH will find architectures that perform well in the restricted low-rank space of their weights, rather than architectures that perform well in the normal training setups, where the weights are no longer restricted. Meanwhile, \enas~allows the weights of its child models to be arbitrary, effectively avoiding such restriction. We suspect this is the reason behind \enas's superior empirical performance to SMASH. In addition, it can be seen from our experiments that \enas~can be flexibly applied to multiple search spaces and disparate domains, \eg~the space of RNN cells for the text domain, the macro search space of entire networks, and the micro search space of convolutional cells for the image domain.
